% Abhakseite „Das kann ich“ – Zone nur im Gesamt (\mitzone)
%% TODO Ich-kann-Sätze: Platzhalter wie in den Aufgabendateien
\begin{abhakseite}
\ifdefined\mitzone
\abhakgruppe{Kennst du schon}
\abhak{1}{Lineare Gleichungssysteme mit zwei und drei Unbekannten lösen, auch unterbestimmte (Parameter frei wählen) und überbestimmte (überzählige Gleichung als Probe)}
\abhak{2}{Vektoren als Listen lesen und schreiben (Zustands-, Mengen-, Kostenvektoren), Komponenten adressieren}
\abhak{3}{Diagramme lesen und zeichnen (Knoten, Pfeile, Punktdiagramme, Kurvenformen unterscheiden)}
\abhak{4}{Terme mit einer Unbekannten aufstellen und gegen Schranken auswerten (linear, mit Randwerten)}
\abhak{5}{Anteile, Prozentsätze und Anzahlen ineinander umrechnen (Bleibe- und Wechselanteile, Gegenanteil)}
\abhak{6}{Potenzen mit Basis unter eins fallen, Exponentialungleichungen durch Probieren oder Logarithmieren lösen, Ergebnisse sinnvoll aufrunden}
\abhak{7}{Lineare Gleichungssysteme mit zwei und drei Unbekannten lösen, auch unterbestimmte (Parameter frei wählen) und überbestimmte (überzählige Gleichung als Probe) – Fehler finden}
\abhak{8}{Lineare Gleichungssysteme mit zwei und drei Unbekannten lösen, auch unterbestimmte (Parameter frei wählen) und überbestimmte (überzählige Gleichung als Probe) – selbst rechnen}
\fi
\abhakgruppe{Einheit 1 · Matrizen als Rechenobjekte}
\abhak{9}{Matrizen als Rechenobjekte}
\abhak{10}{Matrizen als Rechenobjekte – Fehler finden, Begründen}
\abhakgruppe{Einheit 2 · Vektoren unter Matrizen}
\abhak{11}{Vektoren unter Matrizen}
\abhak{12}{Vektoren unter Matrizen – Fehler finden, Begründen}
\abhakgruppe{Einheit 3 · Verflechtung}
\abhak{13}{Verflechtung}
\abhak{14}{Verflechtung – Fehler finden, Begründen, Anwendung}
\abhakgruppe{Einheit 4 · Übergangsmodell}
\abhak{15}{Übergangsmodell}
\abhak{16}{Übergangsmodell – Fehler finden, Begründen, Anwendung}
\abhakgruppe{Einheit 5 · Stationär und langfristig}
\abhak{17}{Stationär und langfristig}
\abhak{18}{Übergangsprozess: Eignung eines Populationsmodells zur langfristigen Beschreibung beurteilen}
\abhak{19}{Stationär und langfristig – Fehler finden, Begründen, Anwendung}
\end{abhakseite}
